\subsection{Dehn 扭转}\label{sec:16c}
令 $M$ 为带角正合辛流形。对每个带标架正合 Lagrangian 球面 $V\subset M$，可赋予其 \emph{Dehn 扭转}，亦称 Picard–Lefschetz 变换，$\tau_V\in\operatorname{Aut}(M,\partial M)$；它在同痕意义下唯一。更准确地说，构造 $\tau_V$ 时作出的附加选择所成的空间可缩。定义时利用 Lagrangian 管状邻域定理，将标架 $v:S\to V$ 延伸为辛嵌入 $\psi:B_s^*S\to\operatorname{int}(M)$，其中 $B_s^*S$ 为 $S$ 上长度 $\leq s$ 的余切向量所成子集，$s>0$。现在取函数 $b:\R\to\R$，满足
\begin{equation}\label{eq:16.5}
\begin{cases}
b(r)=0,&r\geq s,\\
b(r)-b(-r)=-r,&\text{处处成立}.
\end{cases}
\end{equation}
\SourcePage{232}
在 $\psi(B_s^*S)\setminus V$ 上定义 $H(x)=b(|\psi^{-1}(x)|)$，并在 $M\setminus V$ 的其余部分延伸为零。则 $\tau_V$ 是该 Hamilton 函数的时间 $2\pi$ 流，并以唯一可能的方式延伸到 $V$ 本身。为获得更具体的图景，回顾 $T^*S\setminus S$ 上长度函数 $|\xi|$ 的 Hamilton 流可认同为归一化测地流：将任意非零切向量沿其出发的测地线以单位速度移动，而不考虑该向量的长度。因此，$b(|\xi|)$ 的流以速度 $b'(|\xi|)$ 移动切向量。当长度趋于零时，导数趋于 $-1/2$，故极限在零截面上给出对径映射。要证明零截面上的延伸光滑，而不只是连续，可写成 $b(r)=\beta(r^2)-r/2$，其中 $\beta$ 光滑；由于 $b(r)+r/2$ 为偶函数，这是可能的。$\beta(|\xi|^2)$ 的流在整个余切丛上光滑，而 $-|\xi|^2/2$ 的时间 $2\pi$ 流成为对径映射。此外，两 Hamilton 函数均为 $|\xi|$ 的函数，因此 Poisson 交换。由此可将 $b(|\xi|)$ 的时间 $2\pi$ 映射写成两个辛自同构之复合，每一个都显然在零截面附近光滑。\znote{224}{原书此处印作 $|\xi|^2/2$，但它与同一句的分解 $b(r)=\beta(r^2)-r/2$ 以及“导数趋于 $-1/2$”都不相容：$|\xi|^2/2$ 的流是测地流的反向重参数化，单位球面上测地线的周期为 $2\pi$，故时间 $2\pi$ 后回到原处，得到恒等而非对径映射。给出对径映射的应是 $|\xi|/2$：速度为 $-1/2$，时间 $2\pi$ 走弧长 $-\pi$。因此这一句应按 $|\xi|/2$ 理解，$|\xi|^2/2$ 疑为原书笔误。此外，真正在零截面上光滑的只是 $\beta(|\xi|^2)$ 这一支；$-|\xi|/2$ 在零截面上不可微，但由偶性拆分可知其时间 $2\pi$ 流光滑延拓，这正是本段要说明的。[来源：原书 §(16c)，印刷页 224；周期与导数均为初等核对。]}更仔细考察这一构造，可见 $\tau_V$ 实际上正合；这只有在最低维数，即 $M$ 为曲面时才非平凡，此时也可直接观察 Dehn 扭转的标准图形来证明。由定义，类 $[\tau_V]\in\pi_0(\operatorname{Aut}(M,\partial M))$ 在 $V$ 作为带标架正合 Lagrangian 球面的同痕下不变。特别地，它不依赖于 $V$ 的定向。

\StatementNumber{6}\begin{note}\label{st:16.6}
$[\tau_V]$ 是否依赖于标架，是一个颇为微妙的问题。对它诱导的 Fukaya 范畴自等价，已知标架在 $\Ai$ 函子同构意义下不影响结果。这里不证明这一断言，但一个仅涉及对象上作用的较弱版本，可由\stref{推论}{17.17}推出。
\end{note}

取正合 Lefschetz 纤维化 $\pi:E\to S$ 及消失路径 $\gamma$。令 $V_\gamma$ 为相应消失圈，位于 $M=E_{\gamma(0)}$。在 $S\setminus\operatorname{Critv}(\pi)$ 中取绕 $\gamma$ 去而复返的回路 $\lambda$，它逆时针绕过 $\gamma(1)$，见\sr{fig:16.1}{图~16.1}。则沿 $\lambda$ 的单值变换实质上就是沿 $V_\gamma$ 的 Dehn 扭转；更准确地，
\begin{equation}\label{eq:16.6}
h_\lambda\simeq\tau_{V_\gamma},
\end{equation}
其中 $\simeq$ 表示 $\operatorname{Aut}(M)$ 中的同痕。若按\stref{笔记}{15.2}假设纤维化在 $\partial^hE$ 附近平凡，则可稍加强这一断言：同痕可在群 $\operatorname{Aut}(M,\partial M)$ 内进行。这就是\emph{辛 Picard–Lefschetz 定理}。这里不给出证明，但也许值得概述基本策略。首先注意，按\secref{15b}所述作沿 $\partial^hE$ 行为良好的 Lefschetz 纤维化结构变形，不影响结论的有效性，因为它只在同痕意义下改变等式两边。如\secref{15d}所述，可以找到这种变形，使
\SourcePage{233}
$E$ 上辛结构在临界点附近的局部全纯 Morse 坐标中标准化。在这些局部坐标中，以\stref{例题}{16.5}为模型，可显式写出一条短路径趋向临界点时的消失圈。类似地，通过积分有关向量场，可确定绕临界点充分小圆周的单值变换，结果正是沿消失圈的 Dehn 扭转。随后，用一个简单的平行移动论证，就能对原来路径 $\gamma,\lambda$ 推出相应结论。\znote{225}{证明策略可压缩为四步：先用 §(15b) 的沿水平边界良好的变形把辛形式在临界点附近标准化，两边各自只改同痕类，故合法；在标准模型里借例题 16.5 写下显式的短消失路径及其消失圈；对绕临界值的小圆周直接积分向量场，验证单值变换就是沿该消失圈的 Dehn 扭转；最后用平行移动把局部结论搬回任意 $\gamma,\lambda$。要紧的是 $\lambda$ 须在 $\gamma$ 上往返（doubling），这样它恰好绕临界值一周；若绕 $k$ 周则给 $\tau_{V_\gamma}^k$，而不再是同痕意义下的 $\tau_{V_\gamma}$ 本身。[来源：原书 §(16c) 与 §(16d)，印刷页 224–226。]}
\SourceFigure{figure-16-1}{0.1495383}{16.1}

\StatementNumber{7}\begin{note}\label{st:16.7}
显然，用于证明\eref{16.6}的变形并不完全唯一，但可能选择所成的空间可缩。这说明，在两个辛自同构之间，实际上存在唯一的保持端点不动的同痕同伦类。
\end{note}

\clearpage
\subsection{特选基}\label{sec:16d}
现在集中考察底空间为圆盘 $S=D$ 的正合 Lefschetz 纤维化。选择基点 $*\in S\setminus\operatorname{Critv}(\pi)$，并将其视为位于边界上或靠近边界。意思是，若 $*\notin\partial S$，还要选择嵌入的\emph{基路径} $\chi:[0,1]\to S\setminus\operatorname{Critv}(\pi)$，从 $\chi(0)=*$ 通向某点 $\chi(1)\in\partial S$。

消失路径的\emph{特选基}是从 $*$ 出发的消失路径的有序族 $\symbfit\gamma=(\gamma_1,\ldots,\gamma_m)$，$m=|\operatorname{Critv}(\pi)|$，满足下述条件：
\begin{itemize}
\item 不同的 $\gamma_k$ 只在 $*$ 相交，出发方向 $\R^+\gamma_k'(0)\subset TS_*$ 两两不同。
\item 若 $*\in\partial S$，可将 $\R^+\gamma_k'(0)$ 看成上半平面 $H$ 中从原点出发的半直线，故有自然的顺时针顺序。要求它与 $\gamma_k$ 在 $\symbfit\gamma$ 中的顺序一致。另一种情形 $*\notin\partial S$ 下，
\SourcePage{234}
要求每个 $\gamma_k$ 与 $\chi$ 只在 $*$ 相交，且 $\R^+\gamma_k'(0)\ne\R^+\chi'(0)$。此时，$\R^+\gamma_k'(0)$ 是割缝平面 $\CC\setminus\R^+$ 中从原点出发的半直线，其中去掉的方向为 $\chi'(0)$ 的方向。这给出一个顺序，仍施加同样的条件。
\end{itemize}
所有可能的基的同痕类所成集合，带有辫群 $\operatorname{Br}_m$ 通过所谓 \emph{Hurwitz 移动}给出的作用：第 $k$ 个标准生成元 $\beta_k$，$1\leq k<m$，对应于第 $k$ 个初等 Hurwitz 移动 $\symbfit\gamma\mapsto\symbfit\delta=\mathscr L_k\symbfit\gamma$。该移动保持几乎所有路径不变，即 $j\ne k,k+1$ 时 $\delta_j=\gamma_j$，只按\sr{fig:16.2}{图~16.2}变换余下两条路径。容易验证辫关系成立。
\SourceFigure{figure-16-2}{0.44653798}{16.2}

还有一个相关但稍有不同的观点。令 $\mathcal D\subset\operatorname{Diff}^+(S)$ 为这样的保定向微分同胚组成的子群：在 $*$ 上为恒等，适用时在基路径 $\chi$ 上也为恒等，并将子集 $\operatorname{Critv}(\pi)$ 映到自身。显然，$\pi_0(\mathcal D)$ 作用于特选基的同痕类集合；利用 $\operatorname{Diff}(S,\partial S)$ 的可缩性，很容易看出此作用单传递。同类论证还表明，$\pi_0(\mathcal D)$ 同构于 $\operatorname{Br}_m$；然而，此同构并非典范，因此须稍加小心。先回顾：对任意避开 $*$ 及 $\chi$，并满足 $\mu^{-1}(\operatorname{Critv}(\pi))=\{-1,1\}$ 的嵌入路径 $\mu:[-1,1]\to S$，可赋予其\emph{半扭转} $t_\mu\in\mathcal D$，在 $\mu([-1,1])$ 的小邻域之外为恒等，其局部行为由\sr{fig:16.3}{图~16.3}刻画。显然，其中涉及一些选择，但所得类 $[t_\mu]\in\pi_0(\mathcal D)$ 唯一，并在相应路径类别内 $\mu$ 的同痕下不变。给定特选基 $\symbfit\gamma$，可将 $\gamma_k$ 的反向路径与 $\gamma_{k+1}$ 接合，并磨平拐角，定义路径 $\mu_k$，$1\leq k<m$。将标准生成元 $\beta_k$ 映到沿 $\mu_k$ 的半扭转，就给出同构 $\rho_{\symbfit\gamma}:\operatorname{Br}_m\to\pi_0(\mathcal D)$。关于这一同构，初等移动可简写为
\[
\mathscr L_k\symbfit\gamma=t_{\mu_k}^{-1}(\symbfit\gamma)
=\rho_{\symbfit\gamma}(\beta_k^{-1})(\symbfit\gamma).
\]
\SourcePage{235}
\SourceFigure{figure-16-3}{0.28038431}{16.3}
另一方面，这种移动也影响 $\rho_{\symbfit\gamma}$，通过共轭将其变为
\[
\rho_{\mathscr L_k\symbfit\gamma}(\beta)
=t_{\mu_k}^{-1}\circ\rho_{\symbfit\gamma}(\beta)\circ t_{\mu_k}
=\rho_{\symbfit\gamma}(\beta_k^{-1}\beta\beta_k).
\]
结合这两个观察可见，对所有 $k_1,\ldots,k_r\in\{1,\ldots,m-1\}$ 和 $\sigma_1,\ldots,\sigma_r\in\{\pm1\}$，有
\[
\mathscr L_{k_r}^{\sigma_r}\cdots\mathscr L_{k_1}^{\sigma_1}\symbfit\gamma
=\rho_{\symbfit\gamma}(\beta_{k_1}^{-\sigma_1}\cdots\beta_{k_r}^{-\sigma_r})(\symbfit\gamma).
\]
用文字说，以一系列初等移动定义的元素 $\beta\in\operatorname{Br}_m$ 对特选基 $\symbfit\gamma$ 的作用，等于微分同胚 $\rho_{\symbfit\gamma}(\beta^{-1})$ 将 $\symbfit\gamma$ 映到的像。这特别表明，任意两个基都由 Hurwitz 移动联系。\znote{227}{“特选”二字的含义在此：基点固定后，消失路径的排列顺序不是任意的，须与出发方向在 $*$ 处的顺时针顺序一致；所有基的同痕类因此构成一个辫群作用下的单一轨道，$\operatorname{Br}_m$ 经 Hurwitz 移动传递地作用其上。一般位置式的条件（$\gamma_k$ 只在 $*$ 相交、方向互异、避开 $\chi'(0)$）保证移动可逐个进行而不出现退化。注意 $\rho_{\symbfit\gamma}$ 本身依赖基，换基时要共轭，这正是上面共轭公式的来源，也是“辫单值”这一名称的由来：辫群元素记录的是临界值绕基点转动的相对方式。[来源：原书 §(16d)，印刷页 226–227。]}

现在引入消失圈的相应概念。在固定的带角正合辛流形内，考虑由 $m$ 个带标架正合 Lagrangian 球面组成的有序集族 $\symbfit V=(V_1,\ldots,V_m)$。对所涉及的 Lagrangian 球面施加带标架正合同痕所得的集族，不加区分。在此背景下，初等 Hurwitz 移动是某个 $1\leq k<m$ 所给出的变换 $\symbfit V\mapsto\symbfit W=\mathscr L_k\symbfit V$：对所有 $j\ne k,k+1$ 置 $W_j=V_j$，并置
\begin{equation}\label{eq:16.7}
W_k=\tau_{V_k}(V_{k+1}),\qquad W_{k+1}=V_k.
\end{equation}
如通常一样，一般 Hurwitz 移动是这些变换及其逆的任意复合。与先前话题的联系应该很清楚：任意消失路径特选基 $\symbfit\gamma$ 都确定一个有序集族 $\symbfit V=\symbfit V^{\symbfit\gamma}$，由消失圈 $V_k=V_{\gamma_k}$ 组成；称为相应的\emph{消失圈特选基}。\stref{笔记}{16.4}与\eref{16.6}表明，消失路径的初等 Hurwitz 移动 $\mathscr L_k$ 在同痕意义下诱导特选基的相应移动：
\[
\symbfit V^{\mathscr L_k\symbfit\gamma}\simeq\mathscr L_k\symbfit V^{\symbfit\gamma}.
\]
\SourcePage{236}
将集族 $\symbfit V$ 的 \emph{Hurwitz 等价}定义为由 Hurwitz 移动及同痕生成的关系。用这些术语，所得结论就是：$\symbfit V^{\symbfit\gamma}$ 的 Hurwitz 等价类是 Lefschetz 纤维化的不变量；严格地说，是纤维化连同基点选择，以及必要时的基路径选择，的不变量。在适当意义下，它在\secref{15b}定义的沿水平边界行为良好的变形下也不变。
